\subsection{کالبدشکافی ماژول \lr{template\_attention}}

\subsubsection*{۱. فلسفه طراحی، پارادایم معماری و تمثیل ملموس}
ماژول \lr{template\_attention} به عنوان \textbf{معماری مرجع، کالیبراسیون و شاخص مبنا (Standard Benchmark Baseline)} توسعه یافته تا عملکرد تمام مکانیزم‌های دیگر با ساختار استاندارد سنجیده شود.

\textbf{تمثیل ملموس و مثال کاربردی:}
تمثیل \textbf{استوانه استاندارد کیلوگرم و میله متر از آلیاژ پلاتین-ایریدیوم در موزه بین‌المللی اوزان و مقیاس‌های سور پاریس}: جسمی کاملاً خالص، دست‌نخورده و عاری از هرگونه تغییر که تمام ترازوها و وسایل اندازه‌گیری کارخانجات جهان دقت خود را در مقایسه با آن می‌سنجند.

\subsubsection*{۲. فرمول‌بندی و استنتاج عمیق ریاضی \lr{(Mathematical Derivation)}}
\begin{equation}
\operatorname{Attention}(\mathbf{Q}, \mathbf{K}, \mathbf{V}) = \operatorname{Softmax}\left(\frac{\mathbf{Q}\mathbf{K}^T}{\sqrt{d_k}}\right) \mathbf{V}
\end{equation}
بدون اعمال هیچ‌گونه تقریب کم‌مرتبه، هرس تنک یا نگاشت بازگشتی، جهت ارائه بالاترین حد وفاداری عددی ضرب داخلی.

\subsubsection*{۳. کالبدشکافی الگوریتمی و جریان داده \lr{(Algorithmic Tensor Trace)}}
در پیاده‌سازی سورس پایتون در فایل ماژول \lr{template\_attention.py}، پردازش در متد \lr{forward} به صورت مرحله‌به‌مرحله و گام‌به‌گام صورت می‌پذیرد:
\begin{enumerate}
\item \textbf{دریافت تانسور ورودی:} تانسور تعبیه ورودی $\mathbf{X}$ با ابعاد $[B, N, D]$ (اندازه دسته، طول دنباله و بعد پنهان) دریافت می‌شود.
\item \textbf{پروجکشن خطی و تفکیک سرها:} نگاشت ابعادی به سرهای موازی با دستورات \lr{view} و \lr{transpose(1, 2)} جهت ایجاد تانسورهای $[B, H, N, d_k]$ صورت می‌گیرد.
\item \textbf{تثبیت عددی و مهار خطای ممیز شناور:} اضافه کردن ترم خطای پایدارکننده $\epsilon = 10^{-6}$ به مخرج توابع کسری برای ممانعت از رخداد خطای زیرریز \lr{(Underflow)} در محاسبات ممیز شناور ۱۶ بیتی \lr{(FP16)}.
\item \textbf{کاهش ابعاد و تجمیع اطلاعات:} محاسبه ضرب ماتریسی استاندارد کلید و پرس‌وجو، نرمال‌سازی با عملگر سافت‌مکس و ضرب در ماتریس مقدار.
\item \textbf{بازگشت به بعد اصلی و پروجکشن پایانی:} تانسور به شکل $[B, N, D]$ بازآرایی شده و با یک لایه خطی پایانی به لایه پسین تزریق می‌گردد.
\end{enumerate}

\subsubsection*{۴. تحلیل همجوشی سخت‌افزاری و ترافیک حافظه \lr{(Hardware \& Roofline Analysis)}}
از دیدگاه مدل سقف کارایی \lr{(Roofline Model)} و معماری پردازنده‌های گرافیکی مدرن (نظیر کارت گرافیک \lr{NVIDIA RTX 4070}):
\begin{itemize}
\item \textbf{مهار عطش پهنای باند حافظه } \lr{(Memory-Bound Relief)}\rl{:}  در مدل‌های استاندارد، تبادل دائم داده‌های ماتریس توجه با حافظه سراسری کند \lr{DRAM/HBM} باعث معطلی هسته‌های محاسباتی می‌شود. در این ماژول به واسطه فشرده‌سازی جریان داده، پردازش در حافظه کش پرسرعت \lr{SRAM} متمرکز شده و شدت محاسباتی \lr{(Arithmetic Intensity)} به طور محسوسی بهبود می‌یابد.
\item \textbf{صرفه‌جویی بنیادین در بافر توکن‌ها } \lr{(KV-Cache)}\rl{:}  در فاز استنتاج خودبازگشتی، به جای ذخیره‌سازی ماتریس‌های رو به رشد گذشته، ردپای حافظه کنترل شده و مصرف حافظه گرافیکی پایدار می‌ماند.
\end{itemize}

\subsubsection*{۵. ارزیابی تجربی، بنچمارک و توجیه علمی نتایج}
\begin{table}[H]
\centering
\small
\caption{کارنامه تجربی و ارزیابی بنچ‌مارک سخت‌افزاری مکانیزم \lr{template\_attention} روی سه مدل پایه}
\label{tab:benchmark_template_attention}
\begin{tabular}{lccccc}
\toprule
\textbf{مدل پایه \lr{(Base Model)}} & \textbf{دقت \lr{(Accuracy)}} & \textbf{شاخص \lr{F1}} & \textbf{حافظه \lr{(Peak VRAM)}} & \textbf{توان \lr{(Throughput)}} & \textbf{وضعیت} \\
\midrule
\lr{DistilBERT} & \lr{89.91\%} & \lr{0.903} & \lr{329.6 MB} & \lr{71,450 tok/s} & \lr{PASS} \\
\lr{RoBERTa-base} & \lr{58.37\%} & \lr{0.686} & \lr{377.7 MB} & \lr{41,091 tok/s} & \lr{PASS} \\
\lr{GPT-2} & \lr{51.26\%} & \lr{0.676} & \lr{569.8 MB} & \lr{35,033 tok/s} & \lr{PASS\_LOW} \\
\bottomrule
\end{tabular}
\end{table}

\textbf{توجیه علمی و علت‌شناسی جذاب نتایج:}
دقت ۸۹.۹۱٪ و F1 معادل ۹۰.۳۳٪ با مصرف \lr{۳۲۹.۶ MB} در پیاده‌سازی بهینه‌سازی‌شده پایتورک، مقیاس طلایی مقایسه در خط لوله آزمون پایان‌نامه است که نشان می‌دهد ماژول‌های فشرده‌سازی تا چه اندازه توانسته‌اند بدون افت دقت نسبت به این مبنا، حافظه را مهار کنند.

\subsubsection*{۶. دیاگرام معماری شکیل در لاتک \lr{(TikZ Architecture Flowchart)}}
جهت تبیین دقیق جریان داده تانسوری و نحوه تغییر شکل ماتریس‌ها، نمودار معماری این ماژول با استفاده از بسته قدرتمند \lr{TikZ} در محیط لاتک ترسیم شده است:

\begin{center}
\begin{tikzpicture}[scale=0.82, every node/.style={scale=0.82}, 
    node distance=1.1cm,
    auto,
    block/.style={rectangle, draw=blue!70, fill=blue!12, rounded corners=4pt, text centered, minimum height=0.85cm, font=\small},
    kernel/.style={rectangle, draw=orange!80, fill=orange!18, rounded corners=4pt, text centered, minimum height=0.85cm, font=\small},
    outblock/.style={rectangle, draw=teal!80, fill=teal!15, rounded corners=4pt, text centered, minimum height=0.85cm, font=\small}
]
    \node [block] (input) {ورودی $\mathbf{X} \in \mathbb{R}^{B \times N \times D}$};
    \node [kernel, right=1.2cm of input] (k) {نگاشت کرنل $\phi(\mathbf{Q}), \phi(\mathbf{K})$};
    \node [block, right=1.2cm of k] (acc) {تجمع ماتریس حافظه $\mathbf{S}_t$};
    \node [outblock, right=1.2cm of acc] (out) {خروجی خطی};

    \draw [->, thick, draw=blue!80] (input) -- (k);
    \draw [->, thick, draw=orange!80] (k) -- (acc);
    \draw [->, thick, draw=blue!80] (acc) -- (out);
\end{tikzpicture}
\end{center}

\subsubsection*{۷. پیوست تصویری و پرامپت تولید هوش مصنوعی}
برای ترسیم تجسم سه‌بعدی و گرافیک صنعتی این مکانیزم با ابزارهای هوش مصنوعی، پرامپت انگلیسی زیر با دقت تدوین شده است:

\vspace{0.2cm}
\begin{latin}
\footnotesize
\textbf{AI Technical Visualization Prompt:}\\
"A 3D photorealistic render of an indestructible platinum-iridium standard cylinder resting under a double bell jar in a vault, emitting a pristine azure aura."
\end{latin}

\begin{figure}[H]
    \centering
    \framebox[0.9\textwidth]{
    \begin{minipage}{0.85\textwidth}
    \vspace{1.8cm}
    \centering
    \textbf{طرح‌واره معماری و جریان تانسوری ماژول \lr{template\_attention}}\\[0.3cm]\small \lr{(Tensor Flow \& Architectural Schematic for template\_attention)}\\[0.2cm]
    \small (بر اساس پرامپت فنی فوق تولید و در این کادر جایگزین می‌شود)
    \vspace{1.8cm}
    \end{minipage}
    }
    \caption{شماتیک و جریان داده تانسوری در ماژول \lr{template\_attention}}
    \label{fig:template_attention}
\end{figure}

\vspace{0.5cm}
\hrule
\vspace{0.5cm}
